\section{Introduction}

A claims analyst asks a model whether a sentence is supported by a financial
report; a security engineer asks whether a C function contains a vulnerability;
a mathematician asks where the first invalid step occurs. Each case needs a
decision that software can branch on without a second interpretation step.
System One models answer that demand directly. Given a state and a typed
question, Jev returns a binary \noul{}, a candidate \choice{}, or an ordered
\score{} with probabilities \citep{typesafe2026,deusser2026jev}; Laya and Vela
expose similar interfaces in open weights \citep{laya2026,vllm2026vela}.

The interface suggests a simple policy. Ask the cheap binary question first,
read the returned probability, and spend a second call on a differently framed
\choice{} question only when the first answer looks uncertain. The policy is
easy to deploy and matches how practitioners use confidence thresholds. It also
assumes that the probability of the current answer predicts the value of a
different answer. That assumption does not follow from the training objective.
RLCD, the recipe behind Jev and the one Laya documents, rewards honest
probabilities for the question the model answered
\citep{typesafe2026,laya2026,gneiting2007scoring}. A calibrated head can
therefore coexist with a router that has no useful signal. Cascades face the
same limit: confidence-based deferral is optimal only when downstream errors
are tied to upstream confidence, and specialised errors, label noise or
distribution shift can break that link \citep{jitkrittum2023cascade}.

We measure the distinction. Two structured decision
models, Laya 0.4.16 and Vela 2.0 0.3B, are frozen, and routing is studied on
CodeXGLUE defect detection \citep{lu2021codexglue}, FinDVer claim verification
\citep{zhao2024findver} and ProcessBench reasoning-error localisation
\citep{zheng2025processbench}. On the binary tasks the policy pays for \noul{}
and may pay again for \choice{}; on ProcessBench it decides whether to answer or
abstain. We first measure action complementarity, then ask whether the signed
incremental value of \choice{} is predictable from the first call, and finally
compare a genetic algorithm, an evolution strategy and NSGA-II with fixed
policies and budget-matched random search. A one-shot locked test on official
test rows checks the frozen policies.

Gross complementarity is real: on CodeXGLUE, Laya gains 23.3 accuracy points
from a per-row oracle over the better fixed action, and \choice{} repairs 14
\noul{} errors. The same action breaks 15 correct answers, so the net count on
the training pool is $-1$, and held-out pools do not stabilise the sign. A
predictor trained on signed incremental utility has weak out-of-fold rank
correlation, and its confirm cost-adjusted utility is non-positive in every
cell. The genetic algorithm never significantly beats random search; on a
one-shot locked test no genome beats the better fixed action, and the only gain
is over the worse fixed switch, where it learns to stay on \noul{}.
Temperature scaling lowers calibration error but leaves the band policy class
unchanged, so it cannot create the missing signal.

The paper contributes a diagnostic decomposition of routing value into action
complementarity, signed incremental-utility prediction and search efficacy,
measured on frozen structured decision models rather than inferred from
end-task accuracy. It reports a pre-specified negative result across five seeds,
three pools and four task--model cells, and it shows that calibration and
action-value prediction dissociate. The practical lesson is to measure net
complementarity and the predictability of the alternative action before
building a router; a stronger optimiser cannot recover a signal the action set
does not contain.
